\documentclass[11pt]{article}

\usepackage[margin=1in]{geometry}
\usepackage{amsmath,amssymb,bm}
\usepackage{booktabs}
\usepackage{array}
\usepackage{enumitem}
\usepackage[hidelinks]{hyperref}

\title{Detailed Methodology for Condition-Aware Latent\\
State-of-Health Estimation, Forecasting, and\\
Operational Translation in Electric Aircraft Batteries}
\author{
Ben Fogerty \and Chintan Desai \and Shailleze Mahenrenathan \\
Sher Verma \and Sneh Shah \\
\small{MSE 402 Capstone, Group 18}
}
\date{\today}

\begin{document}
\maketitle

\begin{abstract}
This document is the methods companion to the Group 18 capstone Final
Design Review. Its purpose is to explain, with equations and reasoning,
how the project converts noisy battery-management-system (BMS) telemetry
into a usable battery health signal, how that signal is forecast over
short and long horizons, and how it is translated into operational
quantities for an advisory planner. Each section first explains
\emph{why} the design choice was made, then states the equations that
implement it. Reported numerical results match the artifacts written by
the committed pipeline.
\end{abstract}

\tableofcontents

\section{Problem Framing}

The Pipistrel \textsc{velis} Electro reports a battery state of health
(SOH) on every flight and charge event. In principle, this value should
gradually decline as the battery ages. In practice, the reported value
jumps around: it can increase by tens of percentage points between
back-to-back events, behave very differently on charging events than on
flights, and shift whenever the BMS recalibrates its internal capacity
estimate. Treating that raw number as ground truth would mean training a
model to predict noise.

The project therefore separates three quantities:
\begin{itemize}[leftmargin=*]
    \item \(z_k\): the raw BMS-reported SOH at event \(k\). Useful, but
    noisy and context-dependent.
    \item \(x_k\): a latent SOH state, recovered by a Kalman filter that
    treats \(z_k\) as a noisy observation of \(x_k\). This is the signal
    we trust.
    \item \(\widehat{y}_{k,h}\): a forecast of latent SOH \(h\) flights
    ahead of event \(k\), produced by a learned model trained on the
    latent label.
\end{itemize}

Throughout the report, an event is one operational segment for one
battery on one aircraft (a flight, a charge, or an idle gap). Events
are ordered chronologically per battery stream. The committed dataset
contains \(1{,}204\) events: \(1{,}106\) from Plane 166 and \(98\) from
Plane 192, used as a small out-of-distribution holdout.

\section{From Telemetry to Events}

The pipeline starts by reducing each event's raw sample-by-sample
telemetry to a single row of summary features. Two design decisions
matter for everything that follows.

\subsection{Event Type Comes From SOC Direction, Not Labels}

The raw event labels in the data are not always correct. Short ground
operations are sometimes tagged as flights, and a few charge sessions
are tagged ambiguously. Because the downstream noise model treats charge
events differently from flight events, we rebuild the event type from
the data itself.

For each event, we look at how the battery's SOC changed from start to
end. Let \(\Delta\mathrm{SOC}\) be the median SOC over the last samples
of the event minus the median SOC over the first samples. With a small
threshold \(\tau\) to ignore noise,
\begin{equation}
\text{event type} \;=\;
\begin{cases}
\texttt{charge}, & \Delta\mathrm{SOC} \ge \tau, \\
\texttt{flight}, & \Delta\mathrm{SOC} \le -\tau, \\
\texttt{idle},   & |\Delta\mathrm{SOC}| < \tau.
\end{cases}
\end{equation}
A charge event has SOC going up; a flight (or ground discharge) has SOC
going down; an idle gap has SOC essentially flat. This is the only
classification the rest of the pipeline depends on.

\subsection{Per-Event Summary Features}

Within each event, every sample is checked against physical limits (a
plausible current magnitude, voltage range, and temperature range for
that event type). Samples that violate those limits are dropped, then the
event is summarized by a fixed set of statistics:
\begin{itemize}[leftmargin=*]
    \item the median observed SOH \(z_k\),
    \item the SOH inter-quartile range \(\mathrm{IQR}_{\mathrm{SOH},k}\),
    a measure of within-event SOH instability,
    \item mean and 95\textsuperscript{th}-percentile current,
    temperature, and voltage,
    \item SOC minimum, maximum, mean, and span,
    \item the rate-of-change of current and temperature
    (\(|dI/dt|\), \(|dT/dt|\)), summarized as their 95th percentiles,
    \item the gap between the BMS's Kalman SOC estimate and a
    coulomb-counting SOC estimate (this gap grows when the BMS is
    losing track of true SOC).
\end{itemize}
These features are the inputs to both the noise model and the
forecasting models. The reason they matter is that they tell us how
much we should trust the SOH reading at this event.

\section{A Capacity-Based Sanity Check on SOH}

Before introducing the latent model, it is worth showing why we cannot
just compute SOH from physics directly.

On a clean charging event, you can integrate the charging current to get
the total charge delivered, then divide by the SOC range covered to
estimate the battery's full capacity:
\begin{equation}
Q_{\mathrm{delivered}} \;=\; \int |I(t)|\, dt / 3600 \quad (\text{in Ah}),
\qquad
\widehat{C}_{\mathrm{event}} \;=\; \frac{Q_{\mathrm{delivered}}}{(\mathrm{SOC}_{\mathrm{hi}} - \mathrm{SOC}_{\mathrm{lo}})/100}.
\label{eq:coulomb}
\end{equation}
A physics-based SOH is then \(100\,\widehat{C}/C_{\mathrm{rated}}\). The
partial-charging variant of this idea is well-developed in the
literature: Jenu et al.~\cite{jenu2022partial} demonstrate that
integrated-voltage and incremental-capacity methods can give SOH error
below 2\% RMSE on cycled lithium-ion cells under controlled C-rates,
and Saini and Renganathan~\cite{saini2025ica} extend incremental
capacity analysis to dynamic-load electric-vehicle data.

This is a clean idea, but in practice the estimate is noisy on the
aircraft telemetry. On Plane 166, comparing this charge-physics SOH
against the BMS SOH gives \(\mathrm{MAE} \approx 7.9\) SOH points and
\(R^2 \approx 0.11\): the two agree in direction but not closely. We
keep this estimator as a cross-check, but not as the supervised target. A more robust approach
is to treat the BMS reading as a noisy measurement of an underlying
true SOH, and infer the true SOH from many measurements over time.

\section{The Latent SOH Model}

\subsection{Why a Kalman Filter}

Battery aging is gradual: real SOH drifts down slowly over months. The
problem is not that we lack measurements; the problem is that any single
measurement may be wrong by several SOH points because of the operating
context (charging, high current, edge-of-SOC, BMS recalibration). What
we want is a way to combine many noisy measurements over time into a
stable estimate.

This is exactly what a Kalman filter does. It maintains a running
estimate of a hidden quantity, updates that estimate as new measurements
arrive, and weights each measurement by how much we trust it. If a
measurement looks reliable, the estimate moves toward it. If a
measurement looks unreliable, the estimate stays close to its prior.

Treating battery state as a hidden quantity recovered through Kalman
filtering is well established in the battery-management literature. The
foundational treatment by Plett~\cite{plett2004a,plett2004b,plett2004c}
introduced extended Kalman filtering for joint state-of-charge and
state-of-health estimation in lithium-ion battery packs, and subsequent
work has consistently treated SOH as a slowly varying hidden state
inferred from noisy terminal measurements~\cite{yang2021review}. The
mathematical foundations of optimal filtering and smoothing used in this
report follow the standard treatment in
S\"arkk\"a and Svensson~\cite{sarkka2023bayesian}.

\subsection{The State-Space Model}

We model latent SOH as a scalar quantity that drifts slowly over time:
\begin{align}
x_k &= x_{k-1} + w_k, &\quad w_k &\sim \mathcal{N}(0, Q_k) \quad \text{(process)} \label{eq:state_eq}\\
z_k &= x_k + v_k,     &\quad v_k &\sim \mathcal{N}(0, R_k) \quad \text{(observation)} \label{eq:obs_eq}
\end{align}
In words: between events, the latent SOH changes by a small random
amount (Gaussian, mean zero, variance \(Q_k\)). At each event, the BMS
reports a measurement \(z_k\) that equals the true latent SOH plus a
Gaussian observation error with variance \(R_k\).

This is the simplest possible state-space model. We deliberately do not
model degradation rate as a second state, because we do not have enough
events per battery to estimate a rate cleanly. Instead, the rate emerges
from how the latent state drifts over time. The random-walk treatment
of SOH as a slowly drifting hidden parameter is the standard battery-EKF
formulation in Plett~\cite{plett2004c} and is reflected in the recent
review by Yang et al.~\cite{yang2021review}.

\subsection{How Much Drift Is Allowed Between Events: \texorpdfstring{$Q_k$}{Q\_k}}

Real degradation happens on a calendar time scale, not a per-event time
scale. A long idle gap between flights gives the battery more time to
degrade than two back-to-back events. We therefore tie the process
variance to the elapsed time:
\begin{equation}
Q_k \;=\; \sigma_q^{2}\,\Delta t_k,
\label{eq:Q_k}
\end{equation}
where \(\Delta t_k\) is the time between events \(k-1\) and \(k\) in
days and \(\sigma_q\) is the allowed drift per square-root-day. We use
\(\sigma_q = 0.1\) SOH points per \(\sqrt{\text{day}}\). The choice is
deliberately tight: we want the latent state to change slowly, so that
real degradation can drive it down over weeks but a single noisy
measurement cannot push it around.

\subsection{How Much to Trust Each Measurement: \texorpdfstring{$R_k$}{R\_k}}

This is the most important design decision in the project. A fixed
measurement variance would tell the filter to trust every BMS reading
equally. But we know from the data that some readings are far more
reliable than others. A reading taken mid-flight at moderate current
and stable temperature is trustworthy. A reading taken at the top of
charge, right after a BMS recalibration, is not.

The condition-aware approach is to inflate \(R_k\) on events that look
operationally extreme. The filter will then give those events less
influence on the latent state. This idea --- a Kalman filter with
time-varying or adaptive measurement-noise covariance --- has a long
history in the estimation literature, beginning with Mehra's
classification of adaptive filtering approaches~\cite{mehra1972adaptive},
and is the basis for modern adaptive Kalman variants used in lithium-ion
battery state estimation under variable operating
conditions~\cite{yang2021review}. The distinction in our approach is
that, rather than adapting \(R_k\) from innovation residuals after the
fact, we set it \emph{ahead of time} from interpretable physical
condition scores, which makes the noise model transparent and auditable.

We compute several bounded ``stress'' scores for each event, each
between roughly \(0\) (calm) and \(2\) (extreme):
\begin{center}
\begin{tabular}{ll}
\toprule
Score & What it measures \\
\midrule
\(s^{\mathrm{current}}_k\) & How close 95th-percentile current is to the rated maximum \\
\(s^{dI/dt}_k\) & How quickly current was changing during the event \\
\(s^{dT/dt}_k\) & How quickly temperature was changing during the event \\
\(s^{\mathrm{SOCedge}}_k\) & How close the event ran to SOC \(=0\%\) or \(100\%\) \\
\(s^{\mathrm{instab}}_k\) & How variable the SOH reading itself was during the event \\
\(s^{\mathrm{gap}}_k\) & Disagreement between the BMS's Kalman SOC and coulomb-counted SOC \\
\(s^{\mathrm{switch}}_k\) & Whether a BMS recalibration flag fired during the event \\
\(s^{\mathrm{evtype}}_k\) & Larger on charge events than on flights (charges are noisier) \\
\bottomrule
\end{tabular}
\end{center}
The general form of each score is a clipped ratio, for example
\begin{equation}
s^{\mathrm{current}}_k \;=\; \operatorname{clip}\!\left(\frac{\mathrm{p95}|I|_k}{I_{\max}},\; 0,\; 2\right),
\qquad
s^{dI/dt}_k \;=\; \operatorname{clip}\!\left(\frac{\log\!\big(1+\mathrm{p95}|dI/dt|_k\big)}{\log(1+20)},\; 0,\; 2\right).
\end{equation}
The exact upper limits and the use of \(\log(1+\cdot)\) for the rate
scores are just to keep all scores on a comparable scale; another
implementer could pick different normalizations.

A weighted sum of these scores defines a condition multiplier:
\begin{equation}
m_k \;=\; 1 + \sum_{j} w_j\, s^{(j)}_k.
\label{eq:condition_multiplier}
\end{equation}
A baseline ``calm event'' has all scores at zero and \(m_k = 1\). An
extreme event has \(m_k\) several units larger. The weights \(w_j\)
are chosen so that signals we believe more strongly (e.g.\ within-event
SOH variability, BMS recalibration flags) inflate the multiplier more.

Finally, we anchor the multiplier to a per-stream baseline standard
deviation \(\sigma_{\mathrm{base}}\), estimated robustly from
event-to-event SOH changes on the calmest events using the median
absolute deviation. The measurement variance is
\begin{equation}
R_k \;=\; \big(\sigma_{\mathrm{base}}\, m_k\big)^{2},
\qquad
\sigma_{\mathrm{base}} \;\approx\; \frac{1.4826\,\mathrm{MAD}(\Delta z)}{\sqrt{2}}.
\label{eq:R_k}
\end{equation}
The factor \(1.4826\) converts MAD to an equivalent Gaussian standard
deviation; the \(\sqrt{2}\) accounts for the fact that \(\Delta z\) is a
difference of two noisy observations. The result is that
\(R_k\) is small on calm events (the filter trusts them) and large on
extreme events (the filter discounts them).

\subsection{The Kalman Filter Equations}

With latent SOH as a scalar, the filter is especially simple. For each
event \(k = 1, 2, \dots\):
\begin{align}
\widehat{x}_{k|k-1} &= \widehat{x}_{k-1|k-1}             &&\text{(predict the state)} \label{eq:state_pred}\\
P_{k|k-1}           &= P_{k-1|k-1} + Q_k                  &&\text{(predict its variance)} \label{eq:cov_pred}\\
K_k                 &= \frac{P_{k|k-1}}{P_{k|k-1} + R_k}  &&\text{(Kalman gain)} \label{eq:gain}\\
\widehat{x}_{k|k}   &= \widehat{x}_{k|k-1} + K_k\,(z_k - \widehat{x}_{k|k-1}) &&\text{(update with measurement)} \label{eq:update}\\
P_{k|k}             &= (1 - K_k)\,P_{k|k-1}               &&\text{(update its variance)} \label{eq:cov_update}
\end{align}
The Kalman gain \(K_k\) is the key quantity. It is a number between
\(0\) and \(1\). When \(R_k \gg P_{k|k-1}\) the gain is near \(0\) and
the filter essentially ignores the measurement. When \(R_k \ll P_{k|k-1}\)
the gain is near \(1\) and the filter moves the estimate toward the
measurement. The condition-aware \(R_k\) is what controls this
balance event by event.

The output \(\widehat{x}_{k|k}\) is the \emph{causal} latent SOH at
event \(k\), based only on past and current measurements. This is what
we use as the supervised target for forecasting models, because future
information cannot leak into it.

\subsection{The Smoother (for Reports, Not for Training)}

After the forward filter has passed through the entire history, we can
also run a backward pass that revisits each event with the benefit of
future evidence. The standard Rauch--Tung--Striebel (RTS) recursion,
originally derived for maximum-likelihood estimation in linear dynamic
systems~\cite{rauch1965rts}, is
\begin{align}
G_k &= \frac{P_{k|k}}{P_{k+1|k}}, \\
\widehat{x}_{k|N} &= \widehat{x}_{k|k} + G_k\,\big(\widehat{x}_{k+1|N} - \widehat{x}_{k+1|k}\big), \\
P_{k|N} &= P_{k|k} + G_k^{2}\big(P_{k+1|N} - P_{k+1|k}\big).
\end{align}
The smoothed estimate \(\widehat{x}_{k|N}\) is the best retrospective
look at the latent trajectory: it uses every event, not just events up
to \(k\). It is what we plot when we want to visualize how SOH actually
evolved. We never use it as a training target, because it knows about
the future and would leak that information into the forecaster.

\subsection{Result}

The model achieves the central goal: it removes the unphysical jumps in
the raw BMS signal while preserving the long-run trend. On Plane 166:
\begin{center}
\begin{tabular}{lcc}
\toprule
Metric & Raw SOH & Latent SOH \\
\midrule
Total event-to-event variation (Battery 1) & 484 & 47 \\
Total event-to-event variation (Battery 2) & 469 & 45 \\
Largest single upward jump (Battery 1) & 29 & 0.7 \\
Largest single upward jump (Battery 2) & 25 & 0.7 \\
\bottomrule
\end{tabular}
\end{center}
The raw signal can apparently ``recover'' twenty-nine SOH points in a
single event. The latent signal cannot, by construction, because the
condition-aware filter recognized those jumps as coming from
unreliable measurements.

\section{Forecasting Latent SOH}

Once we have a stable latent SOH, the next question is how it will
evolve. We frame this as: given the latent SOH at event \(k\) and the
event history leading up to it, predict the latent SOH \(h\) flights
into the future.

\subsection{Why a Sequence Model}

Battery degradation depends on history, not just on the current state.
Two batteries with the same SOH today can degrade at very different
rates depending on how they have been used --- frequent fast charges,
high temperatures, deep discharges, idle storage at high SOC, and so on.
A model that only sees today's state cannot tell those two batteries
apart.

Recurrent neural networks (LSTM and GRU) are designed exactly for this:
they consume a sequence of past events and produce a single prediction
that depends on the whole sequence. LSTM-based SOH estimators have been
demonstrated in the battery-health literature with mean-absolute
percentage errors below \(2.22\%\) on capacity-test
data~\cite{yu2022soh}, which motivated our choice of recurrent models
as the primary forecasting family. We use them here to predict SOH
\(h\) flights ahead.

\subsection{Target and Inputs}

For each event \(k\), the target is the latent SOH at the future event
that lands \(h\) flights later:
\begin{equation}
y_{k,h} \;=\; x^{\mathrm{filter}}_{j(k,h)},
\qquad
j(k,h) \;=\; \min\{ j > k \,:\, f_j \ge f_k + h\},
\end{equation}
where \(f_k\) is the cumulative flight count for that battery up to
event \(k\). Using cumulative flight count rather than calendar time
ensures that charge events and idle gaps do not distort the horizon.

The model input at event \(k\) is a window of the last \(L\) events,
each represented by a feature vector \(\mathbf{u}_t\) containing the
summary statistics from Section~2.2 plus engineered features
(cumulative throughput, rolling averages, lagged SOH, etc.). Concretely,
\begin{equation}
\mathbf{X}_k \;=\;
\begin{bmatrix}
\mathbf{u}_{k-L+1} \\
\mathbf{u}_{k-L+2} \\
\vdots \\
\mathbf{u}_k
\end{bmatrix} \in \mathbb{R}^{L\times p}.
\end{equation}
Features are standardized using training-split statistics so that the
network sees comparable scales across columns.

\subsection{Predicting a Change, Not a Level}

A useful trick is to ask the network to predict the \emph{change} in
SOH rather than the next SOH directly:
\begin{equation}
\widehat{y}_{k,h} \;=\; x^{\mathrm{filter}}_k + g_\theta(\mathbf{X}_k).
\label{eq:delta_pred}
\end{equation}
The network output \(g_\theta(\mathbf{X}_k)\) is a small delta, and
adding it back to the current SOH gives the level. This focuses the
network's capacity on the part that actually requires learning ---
incremental degradation --- rather than on relearning the level of SOH
on every event.

\subsection{The Recurrent Cells}

We use LSTM and GRU as the recurrent core. Both consume the lookback
window step by step and produce a final hidden state that summarizes
the sequence. For one step \(t\), the LSTM cell is
\begin{align}
\mathbf{i}_t &= \sigma(\mathbf{W}_i\mathbf{u}_t + \mathbf{U}_i\mathbf{h}_{t-1} + \mathbf{b}_i),
& \text{input gate} \\
\mathbf{f}_t &= \sigma(\mathbf{W}_f\mathbf{u}_t + \mathbf{U}_f\mathbf{h}_{t-1} + \mathbf{b}_f),
& \text{forget gate} \\
\mathbf{o}_t &= \sigma(\mathbf{W}_o\mathbf{u}_t + \mathbf{U}_o\mathbf{h}_{t-1} + \mathbf{b}_o),
& \text{output gate} \\
\widetilde{\mathbf{c}}_t &= \tanh(\mathbf{W}_c\mathbf{u}_t + \mathbf{U}_c\mathbf{h}_{t-1} + \mathbf{b}_c),
& \text{candidate cell} \\
\mathbf{c}_t &= \mathbf{f}_t \odot \mathbf{c}_{t-1} + \mathbf{i}_t \odot \widetilde{\mathbf{c}}_t,
& \text{cell state} \\
\mathbf{h}_t &= \mathbf{o}_t \odot \tanh(\mathbf{c}_t).
& \text{hidden state}
\end{align}
The three sigmoidal gates control how much past information is
retained, how much new information is admitted, and how much of the
internal state is exposed. The final hidden state \(\mathbf{h}_L\) is
passed through a small multilayer perceptron head to produce the scalar
delta prediction.

The GRU is a simpler variant of the same idea, using two gates instead
of three:
\begin{align}
\mathbf{z}_t &= \sigma(\mathbf{W}_z\mathbf{u}_t + \mathbf{U}_z\mathbf{h}_{t-1} + \mathbf{b}_z),
& \text{update gate} \\
\mathbf{r}_t &= \sigma(\mathbf{W}_r\mathbf{u}_t + \mathbf{U}_r\mathbf{h}_{t-1} + \mathbf{b}_r),
& \text{reset gate} \\
\widetilde{\mathbf{h}}_t &= \tanh(\mathbf{W}_h\mathbf{u}_t + \mathbf{U}_h(\mathbf{r}_t \odot \mathbf{h}_{t-1}) + \mathbf{b}_h), \\
\mathbf{h}_t &= (1-\mathbf{z}_t)\odot \mathbf{h}_{t-1} + \mathbf{z}_t \odot \widetilde{\mathbf{h}}_t.
\end{align}
We try both because GRUs are sometimes easier to train on small
datasets while LSTMs sometimes hold longer dependencies more stably.
The winning model is chosen by validation MAE at each horizon.

\subsection{Loss and Optimization}

Training minimizes the Smooth L1 (Huber) loss on the delta target:
\begin{equation}
\mathcal{L}(\theta) \;=\; \frac{1}{N}\sum_{k=1}^{N}\operatorname{SmoothL1}\!\big(g_\theta(\mathbf{X}_k),\; y_{k,h}-x^{\mathrm{filter}}_k\big),
\end{equation}
where
\begin{equation}
\operatorname{SmoothL1}(\widehat{u}, u) =
\begin{cases}
\tfrac{1}{2}(\widehat{u}-u)^{2}, & |\widehat{u}-u| < 1, \\
|\widehat{u}-u| - \tfrac{1}{2}, & \text{otherwise.}
\end{cases}
\end{equation}
Huber loss is used in preference to plain MSE because it is less
sensitive to a small number of large residuals, which is helpful given
the noisiness of the underlying problem. We optimize with Adam,
gradient clipping, and early stopping on validation MAE.

\subsection{Validation: Why Walk-Forward}

A standard random train/test split would let the model peek at events
from the future of its training history, which is exactly the kind of
leakage we are trying to avoid. We use a chronological walk-forward
scheme instead: train on the earliest fraction of events, validate on
the next contiguous chunk, then advance. Plane 192 is kept entirely
out of training and used only to check that the model has not
overfitted to Plane 166.

\subsection{Result}

Forecast quality depends strongly on horizon. The committed best models
and their backtest mean absolute errors are:
\begin{center}
\begin{tabular}{lll}
\toprule
Horizon (flights ahead) & Best model & Backtest MAE (SOH points) \\
\midrule
1 & LSTM & 0.19 \\
5 & LSTM & 0.66 \\
10 & GRU & 1.14 \\
15 & GRU & 1.17 \\
20 & GRU & 1.19 \\
\bottomrule
\end{tabular}
\end{center}
Errors at one and five flights are well within useful precision for
operational planning. Errors at fifteen and twenty flights are about
six times larger than at one flight. This is a meaningful and honest
result: the data supports confident short-horizon prediction, but not
confident long-horizon prediction. The long-horizon problem requires a
different approach.

\section{Physics-Informed Forecasters}

Two of the trained models add physical structure to the recurrent
forecaster. They are useful both for performance and for
interpretability.

\subsection{Why Add Physics}

A pure recurrent network is a flexible function approximator. With
limited data, that flexibility can lead it to fit patterns that are not
physically meaningful: SOH might be predicted to recover, or to drop
faster on a cool flight than on a hot one. Physics-informed structure
constrains the model to behave like a battery: degradation only goes
one way, certain stresses (heat, throughput, time) increase
degradation, and the degradation rate should change smoothly from one
similar event to the next.

We use two such models: a \emph{hybrid} model that hard-codes battery
aging intuition into its architecture, and a \emph{standard PINN} that
enforces a differential equation through automatic differentiation. The
PINN formulation is inspired by Gu et al.~\cite{gu2025ppinn}, who
parameterize a single-particle electrochemical model with a
physics-informed neural network and identify internal aging parameters
from partial charging profiles. Our PINN does not solve the full
electrochemical PDE --- our event-level data does not support that ---
but it borrows the idea of embedding a physical degradation law into the
training loss.

\subsection{Aging Channels}

Both physics-informed models decompose degradation into five physically
named drives, each a nonnegative scalar built from the event summary
features. The intuition behind each:
\begin{itemize}[leftmargin=*]
    \item \(D^{\mathrm{cal}}_k\) (calendar): captures aging that happens
    simply because time passed. Larger on long idle gaps, especially at
    high storage SOC.
    \item \(D^{\mathrm{cyc}}_k\) (cycling): captures aging from
    throughput. Larger on long, high-current discharges with wide SOC
    swings.
    \item \(D^{\mathrm{th}}_k\) (thermal): captures aging from heat.
    Built from an Arrhenius-style temperature factor and time spent
    above warm temperatures.
    \item \(D^{\mathrm{res}}_k\) (resistance/health-fault): captures
    aging signals from voltage sag and from inconsistency between the
    BMS's own SOC estimators.
    \item \(D^{\mathrm{hist}}_k\) (history): captures accumulated
    cumulative stress, including total equivalent full cycles and
    rolling stress.
\end{itemize}
The actual definitions are smooth nonnegative combinations of upstream
features wrapped in \(\log(1 + \cdot)\) to keep them on a comparable
scale. A different team could choose somewhat different formulas; the
important property is that all five drives are nonnegative and that
they correspond to distinct aging mechanisms.

\subsection{The Hybrid Model}

The hybrid model uses the aging channels directly in its prediction
rule. A small neural network looks at the event context and produces
five nonnegative sensitivities \(\alpha_{k,1}, \dots, \alpha_{k,5}\),
one per channel. The base degradation at event \(k\) is then the
dot product
\begin{equation}
\mathcal{D}^{\mathrm{base}}_k \;=\; \sum_{j=1}^{5} \alpha_{k,j}\, D^{(j)}_k.
\end{equation}
Two additional terms refine this: a small interaction term that lets
channels reinforce one another (high temperature \emph{and} high
throughput is worse than either alone), and a learned residual term
that captures whatever the explicit channels miss. The total
degradation is subtracted from the current SOH to give the prediction:
\begin{equation}
\widehat{x}_{k+1} \;=\; \operatorname{clip}\!\big(x^{\mathrm{filter}}_k - \mathcal{D}_k,\; 0,\; 100\big).
\label{eq:hybrid_pred}
\end{equation}
Two structural properties matter:
\begin{itemize}[leftmargin=*]
    \item The sensitivities \(\alpha_{k,j}\) are forced to be
    nonnegative (via a softplus), so degradation cannot become negative
    and SOH cannot ``recover'' through the explicit channels.
    \item The sensitivities are encouraged to vary \emph{smoothly}
    between nearby events on the same battery. The training loss
    includes a penalty on \(\|\alpha_k - \alpha_{k-1}\|^2\), weighted
    by an exponential decay over the inter-event gap. This means the
    same battery cannot suddenly change its aging law from one event
    to the next.
\end{itemize}

\subsection{The Standard PINN}

The PINN takes a more classical approach. It treats SOH as a continuous
function of battery age \(t\), \(u_\theta(t)\), and asks the network
to satisfy a governing differential equation:
\begin{equation}
\frac{du_\theta}{dt} \;+\; g(\mathbf{D}_k,\, u_\theta) \;=\; 0,
\label{eq:pinn_ode}
\end{equation}
where \(g(\cdot)\) is a nonnegative degradation rate built from the
same aging channels and a small set of trainable weights:
\begin{equation}
g(\mathbf{D}, u) \;=\; \big(b + \mathbf{w}^{\top}\mathbf{D}\big)\,\Big(1 + \gamma\,\tfrac{100-u}{100}\Big),
\qquad b, \mathbf{w}, \gamma \ge 0.
\end{equation}
Reading the equation: the rate of SOH loss is positive, increases with
the aging drives, and increases as the battery moves further from
fresh.

The clever part is how the PINN is trained. We do not need to discretize
the ODE; PyTorch's automatic differentiation gives us
\(du_\theta/dt\) directly. We then add a residual loss that penalizes
the network whenever
\(du_\theta/dt + g(\mathbf{D}, u_\theta)\) departs from zero. The full
training objective combines:
\begin{itemize}[leftmargin=*]
    \item a data loss matching \(u_\theta(t_k)\) to observed SOH,
    \item a step loss matching the Euler-stepped forecast to observed
    next SOH,
    \item the ODE-residual loss above,
    \item a monotonic prior penalizing any positive
    \(du_\theta/dt\) (SOH should not go up over time),
    \item an initial-condition term anchoring the first event.
\end{itemize}
The PINN is more conceptually faithful to the underlying physics than
the hybrid model, but it is harder to train on the available dataset
because the residual loss must be balanced against the other terms.

\section{Long-Horizon Outlook via an External Backbone}

\subsection{Why We Cannot Forecast EOL From Our Data Alone}

The aircraft batteries in the dataset have not yet aged to the end of
their life. We have observed perhaps thirty percent of a typical
battery's lifetime; we have never seen one reach the replacement
threshold. Any model that tries to extrapolate to end-of-life from our
data alone is extrapolating beyond its training range.

The fix is to borrow the \emph{shape} of degradation from an external
dataset that does cover full battery lifetimes, then place our
batteries on that shape.

\subsection{The Backbone Shape}

We use the public eVTOL battery test dataset from Bills et
al.~\cite{bills2023evtol} (Carnegie Mellon University), which contains
detailed cycle-life experiments on lithium-ion cells operated under
representative eVTOL mission profiles. The dataset spans many cells
cycled to end of life under varying cruise lengths, charge currents,
voltage limits, and temperatures, making it well suited for learning a
generic eVTOL degradation shape. For each test battery, we map calendar
progress to normalized progress \(p \in [0,1]\), with \(p = 0\) at the
start of life and \(p = 1\) at end of life. We also map measured capacity to an
adjusted health scale where \(100\) corresponds to fresh and \(0\)
corresponds to the end-of-life capacity threshold.

Pooling all external test points gives a cloud of
\((p_a,\, \mathrm{Health}^{\mathrm{adj}}_a)\) pairs. We fit a
monotone-decreasing function \(B(p)\) through this cloud using isotonic
regression:
\begin{equation}
B(p) \;=\; \arg\min_{\substack{f:[0,1]\to[0,100] \\ f\text{ non-increasing}}}
\sum_{a}\big(\mathrm{Health}^{\mathrm{adj}}_a - f(p_a)\big)^{2}.
\label{eq:isotonic}
\end{equation}
The result is a smooth, decreasing curve that captures the typical
degradation trajectory of an eVTOL battery from new to retired.

\subsection{Placing an Aircraft Battery on the Backbone}

For each of our aircraft batteries, we estimate where on the backbone
the battery currently sits. We grid-search two quantities:
\begin{itemize}[leftmargin=*]
    \item \(p_0\): the normalized progress at which the battery's
    observable life starts (because we do not have data from new).
    \item \(L\): the total life of the battery in flights.
\end{itemize}
For each candidate pair, we predict latent SOH at every observed
flight using the backbone, then minimize the mean absolute error:
\begin{equation}
(p_0^{\star}, L^{\star}) \;=\;
\arg\min_{p_0, L} \frac{1}{n}\sum_{k=1}^{n}\bigg|x^{\mathrm{filter}}_k - B\!\Big(p_0 + \tfrac{f_k - f_1}{L}\Big)\bigg|.
\label{eq:plane_calibration}
\end{equation}
Once the best \((p_0^{\star}, L^{\star})\) is known, the remaining
flights to end of life are simply \((1 - p^{\star}_{\mathrm{end}})\,L^{\star}\),
where \(p^{\star}_{\mathrm{end}}\) is the progress at the most recent
event.

For Plane 166, the fitted MAE between the calibrated backbone and the
latent SOH trajectory is approximately \(1.7\) SOH points across both
batteries, which means the external shape genuinely captures the
long-run behavior of our aircraft batteries even though we have not
observed them anywhere near end of life.

\section{From SOH to Operational Quantities}

The final piece of the pipeline turns SOH into things an operator
actually plans around: how many circuits the aircraft can fly today,
how fast the battery will drain on a given flight, and when the
battery will need to be replaced.

\subsection{Circuits per Flight}

The Pilot Operating Handbook gives a baseline relationship between SOH
and SOC consumed per traffic circuit, in the form of a small lookup
table. We turn this into a continuous function by linear interpolation:
\begin{equation}
s_{\mathrm{POH}}(\mathrm{SOH}) \;=\; \operatorname{interp}\!\big(\mathrm{SOH};\; [0, 20, 40, 60, 80, 100],\; [20, 16, 13, 12, 10, 9]\big).
\end{equation}
The reading: a fresh battery uses about 9\% SOC per circuit; a heavily
degraded battery uses about 20\%.

The POH curve is generic; each aircraft uses slightly more or less SOC
per circuit than the table predicts. We calibrate a per-plane scalar
multiplier \(k^{\star}\) by exploiting a structural fact: real flights
fly an integer number of circuits. For any candidate \(k\), the
inferred number of circuits during event \(k\) is
\begin{equation}
\widehat{n}_k(k) \;=\; \frac{\Delta\mathrm{SOC}_k}{k\,s_{\mathrm{POH}}(\mathrm{SOH}_k)},
\end{equation}
and we pick the \(k\) that makes \(\widehat{n}_k\) cluster most tightly
around integers:
\begin{equation}
k^{\star} \;=\; \arg\min_{k}\;\frac{1}{n}\sum_{k=1}^{n}\Big|\widehat{n}_k(k) - \operatorname{round}\!\big(\widehat{n}_k(k)\big)\Big|.
\end{equation}
The result is \(k^{\star}_{166} = 0.80\) for Plane 166 (slightly more
efficient than the POH baseline) and \(k^{\star}_{192} = 0.74\) for
Plane 192. The number of feasible circuits at start SOC \(S_0\) with a
30\% reserve is then
\begin{equation}
N_{\max} \;=\; \left\lfloor \frac{\max(S_0 - 30,\; 0)}{k^{\star}\,s_{\mathrm{POH}}(\mathrm{SOH})} \right\rfloor.
\end{equation}

\subsection{SOC Burn Rate}

Circuits give a coarse view of capacity. For finer planning, we
predict the SOC burn rate during flights:
\begin{equation}
r_k \;=\; \frac{\mathrm{SOC \;drop}_k}{\mathrm{duration}_k\;\text{(minutes)}}.
\end{equation}
We fit a linear regression
\begin{equation}
\widehat{r}_k \;=\; \beta_0 + \sum_j \beta_j\, x_{k,j}
\end{equation}
on a small set of physically meaningful features: latent SOH, mean and
peak current, mean cell temperature, mean voltage, SOC range, current
rate-of-change, and the BMS coulomb-gap signal. The linear form is
chosen deliberately: SOC burn rate is much more directly observable
than SOH, the relationship is approximately linear in the physically
meaningful features, and we want the model to be interpretable.

The fit is strong. On the Plane 166 test split, the model achieves
\(\mathrm{MAE} = 0.19\) SOC\%/min and \(R^2 = 0.56\); on Plane 192 (out
of distribution) it achieves \(R^2 = 0.54\). This is the most directly
useful piece of the operational layer, because mission energy planning
depends more on SOC burn than on SOH per se.

\section{Putting It All Together}

End to end, the system performs the following:

\begin{enumerate}[leftmargin=*]
    \item Parse raw telemetry into events. Reclassify event type from
    SOC direction so charge and flight events are correctly identified.
    Summarize each event into a single feature row.
    \item Detect how trustworthy each event's SOH measurement is via
    bounded condition scores. Combine them into an event-specific
    measurement variance \(R_k\).
    \item Run a Kalman filter that treats latent SOH as a slowly
    drifting hidden state and observed BMS SOH as a noisy measurement
    of it. The filter automatically down-weights extreme events
    because their \(R_k\) is larger.
    \item Use the causal filtered latent SOH as the supervised target
    for forecasting models. Train LSTM and GRU sequence models in
    delta mode with Huber loss and walk-forward validation. Use
    physics-informed variants (hybrid, PINN) where physical structure
    helps.
    \item For long-horizon outlook, calibrate an isotonic backbone
    learned from an external eVTOL dataset that does see end of life,
    and place each of our batteries on that backbone via a two-parameter
    grid search.
    \item Translate SOH into operational quantities: feasible
    circuits per flight (from a POH-derived map with per-plane
    calibration) and SOC burn rate (from a small linear regression on
    physically meaningful features).
\end{enumerate}

\section{Limitations}

\begin{itemize}[leftmargin=*]
    \item Latent SOH is modeled as a scalar random walk. A two-state
    model with explicit degradation rate would be more expressive but
    requires more data per battery than we have.
    \item The measurement variance \(R_k\) is built from interpretable
    but heuristic stress scores. The weights were validated by checking
    that the filter residuals were well-behaved, not derived from a
    first-principles noise model.
    \item Short-horizon forecasts (one to five flights) are reliable.
    Beyond about ten flights, forecast error grows substantially, which
    is why the long-horizon outlook uses backbone calibration rather
    than the recurrent forecaster.
    \item Plane 192 has only 98 events; it is a useful sanity check but
    is too small to be a strong claim of cross-aircraft generalization.
    \item The PINN enforces an ODE, not a partial differential equation.
    A full electrochemical PINN (single-particle or Doyle--Fuller--Newman)
    would require dense voltage--current trajectories that we do not
    have.
\end{itemize}

\section{Why Each Choice Was Made}

Distilled to a single paragraph: raw BMS SOH was rejected as a label
because empirically it jumps by tens of points in a single event. A
Kalman filter was chosen because it is the simplest principled way to
combine many noisy measurements of a slowly-drifting hidden quantity.
The filter is condition-aware because we know which events are
unreliable, and ignoring that knowledge would waste it. Sequence models
were chosen because aging depends on history, not just on present
state. Physics-informed variants were added to constrain learning when
data are limited. An external backbone was added because our aircraft
data does not span a full battery lifetime. The operational layer was
kept simple --- linear models and lookup tables --- because the
quantities it predicts are directly observable and a complex model
would not gain accuracy. Throughout, the goal was not to maximize
model complexity but to make each step as transparent as the data
would support.

\begin{thebibliography}{9}

\bibitem{plett2004a}
G.~L. Plett,
``Extended Kalman filtering for battery management systems of
LiPB-based HEV battery packs: Part 1. Background,''
\emph{Journal of Power Sources}, vol.~134, no.~2, pp.~252--261, 2004.
doi: 10.1016/j.jpowsour.2004.02.031.
\emph{(Foundational treatment of EKF for battery management. Establishes
the conceptual frame we follow: battery state quantities are hidden,
estimated from noisy terminal measurements.)}

\bibitem{plett2004b}
G.~L. Plett,
``Extended Kalman filtering for battery management systems of
LiPB-based HEV battery packs: Part 2. Modeling and identification,''
\emph{Journal of Power Sources}, vol.~134, no.~2, pp.~262--276, 2004.
doi: 10.1016/j.jpowsour.2004.02.032.
\emph{(State-space modeling and parameter identification for lithium
battery packs in the Kalman framework.)}

\bibitem{plett2004c}
G.~L. Plett,
``Extended Kalman filtering for battery management systems of
LiPB-based HEV battery packs: Part 3. State and parameter estimation,''
\emph{Journal of Power Sources}, vol.~134, no.~2, pp.~277--292, 2004.
doi: 10.1016/j.jpowsour.2004.02.033.
\emph{(Treats SOH as a slowly varying parameter recovered through
recursive Kalman estimation. Direct precedent for our random-walk
latent-SOH formulation.)}

\bibitem{rauch1965rts}
H.~E. Rauch, F.~Tung, and C.~T. Striebel,
``Maximum likelihood estimates of linear dynamic systems,''
\emph{AIAA Journal}, vol.~3, no.~8, pp.~1445--1450, 1965.
doi: 10.2514/3.3166.
\emph{(Original paper deriving the backward-pass smoother we use for
retrospective latent-SOH estimation.)}

\bibitem{mehra1972adaptive}
R.~K. Mehra,
``Approaches to adaptive filtering,''
\emph{IEEE Transactions on Automatic Control}, vol.~AC-17, no.~5,
pp.~693--698, 1972.
doi: 10.1109/TAC.1972.1100100.
\emph{(Classical reference for Kalman filtering with time-varying noise
covariance. Establishes the legitimacy of adapting \(R_k\), which our
condition-aware noise model does in an a-priori, physics-driven way
rather than from innovation statistics.)}

\bibitem{sarkka2023bayesian}
S.~S\"arkk\"a and L.~Svensson,
\emph{Bayesian Filtering and Smoothing}, 2nd ed.
Cambridge University Press, 2023.
\emph{(Standard modern textbook for Kalman filtering and RTS smoothing.
Used as the methodological reference for the linear-Gaussian formulation
and notation in this report.)}

\bibitem{yang2021review}
S.~Yang, C.~Zhang, J.~Jiang, W.~Zhang, L.~Zhang, and Y.~Wang,
``Review on state-of-health of lithium-ion batteries: Characterizations,
estimations and applications,''
\emph{Journal of Cleaner Production}, vol.~314, p.~128015, 2021.
doi: 10.1016/j.jclepro.2021.128015.
\emph{(Recent review surveying SOH estimation techniques. Confirms that
adaptive Kalman variants with condition-dependent noise are an
established branch of the SOH-estimation literature, consistent with our
approach.)}

\bibitem{yu2022soh}
Z.~Yu, Y.~Zhang, L.~Qi, and R.~Li,
``SOH Estimation Method for Lithium-ion Battery Based on Discharge
Characteristics,''
\emph{International Journal of Electrochemical Science},
vol.~17, no.~7, Article 220725, 2022.
doi: 10.20964/2022.07.38.
\emph{(Uses an LSTM-based degradation model for battery SOH estimation
under partial charging/discharging; supports our choice of recurrent
sequence models for short-horizon forecasting.)}

\bibitem{jenu2022partial}
S.~Jenu, A.~Hentunen, J.~Haavisto, and M.~Pihlatie,
``State of health estimation of cycle aged large format lithium-ion
cells based on partial charging,''
\emph{Journal of Energy Storage}, vol.~46, p.~103855, 2022.
doi: 10.1016/j.est.2021.103855.
\emph{(Develops incremental-capacity and integrated-voltage SOH
estimation from partial charging segments. Background for our charge-event
coulomb-counting cross-check estimator.)}

\bibitem{saini2025ica}
U.~Saini and S.~Renganathan,
``Incremental capacity analysis of battery under dynamic load
conditions,''
\emph{MethodsX}, vol.~14, p.~103331, 2025.
\emph{(Extends ICA-style SOH estimation to electric-vehicle batteries
under non-laboratory dynamic loads. Motivates our use of charge-event
SOH only on clean, monotonic charge segments and our decision to treat
raw BMS SOH as a noisy observation rather than ground truth.)}

\bibitem{gu2025ppinn}
X.~Gu, X.~Huan, Y.~Ren, W.~Zhou, W.~Jiang, and Z.~Song,
``Real-Time Physics-Aware Battery Health Monitoring from Partial Charging
Profiles via Physics-Informed Neural Networks,''
arXiv preprint arXiv:2511.12053, 2025.
\emph{(Parameterizes a single-particle electrochemical model with a
PINN to identify internal aging parameters from partial charging.
Methodological inspiration for our physics-hybrid and ODE-form PINN
forecasters, which adapt the idea of physics-informed losses to
event-level aircraft telemetry rather than dense charge curves.)}

\bibitem{bills2023evtol}
A.~Bills, S.~Sripad, L.~Fredericks, M.~Guttenberg, D.~Charles, E.~Frank,
and V.~Viswanathan,
``eVTOL Battery Dataset,''
Carnegie Mellon University, Pittsburgh, PA, USA.
\emph{(Public lithium-ion cycle-life dataset under representative eVTOL
mission profiles. Provides the external long-life degradation shape
that our isotonic backbone is fit to, enabling long-horizon remaining-life
calibration despite our aircraft batteries not yet having aged to end of
life.)}

\end{thebibliography}

\end{document}
